\begin{solution}{95}
\begin{subsolution}
设该原子核从静止的第一激发态 $^{57}Fe^{*}$ 跃迁到基态 $^{57}Fe$ ，发出一频率为 $\nu_{e}$ 的光子，同时基态 $^{57}Fe$ 获得一反冲速度 $v$（发射过程）。由能量、动量守恒有
\begin{equation}
E _ {0} = h \nu_ {\mathrm{e}} + \frac {1}{2} m _ {\mathrm{Fe}} v ^ {2}
\label{eq:1.s95}
\end{equation}
\begin{equation}
0 = \frac {h \nu_ {\mathrm{e}}}{c} - m _ {\mathrm{Fe}} v
\label{eq:2.s95}
\end{equation}
将\eqref{eq:2.s95}式代入\eqref{eq:1.s95}式得，在考虑核反冲的情形下，从激发态跃迁到基态发出的光子的频率为
\begin{equation}
\nu_ {\mathrm{e}} = \frac {E _ {0}}{h} - \frac {h \nu_ {\mathrm{e}} ^ {2}}{2 m _ {\mathrm{Fe}} c ^ {2}} = \nu_ {0} - \frac {h \nu_ {\mathrm{e}} ^ {2}}{2 m _ {\mathrm{Fe}} c ^ {2}}
\label{eq:3.s95}
\end{equation}
式中
\begin{equation}
\nu_ {0} = \frac {E _ {0}}{h}
\label{eq:4.s95}
\end{equation}
是在不考虑核反冲情形下 $^{57}$Fe 从第一激发态跃迁到基态发出的光子的频率，或从基态跃迁到第一激发态吸收的光子的频率。

注意到 $\frac{h\nu_{e}}{m_{Fe}c^{2}}<<1$ ，\eqref{eq:3.s95}式右边的 $\nu_{e}$ 可以用 $\nu_{0}$ 代替，于是有
\begin{equation}
\nu_ {\mathrm{e}} \approx \nu_ {0} - \frac {h \nu_ {0} ^ {2}}{2 m _ {\mathrm{Fe}} c ^ {2}} = \nu_ {0} \left(1 - \frac {E _ {0}}{2 m _ {\mathrm{Fe}} c ^ {2}}\right)
\label{eq:5.s95}
\end{equation}
在考虑核反冲和不考虑核反冲的情形下,从第一激发态跃迁到基态发出的光子的频率之差为
\begin{equation}
\nu_ {\mathrm{e}} - \nu_ {0} \approx - \nu_ {0} \frac {E _ {0}}{2 m _ {\mathrm{Fe}} c ^ {2}} = - 4. 7 \times 1 0 ^ {1 1} \mathrm{Hz}
\label{eq:6.s95}
\end{equation}
由\eqref{eq:1.s95}\eqref{eq:6.s95}式得，反冲能 $E_{\mathrm{R}}=\frac{1}{2}m_{\mathrm{Fe}}v^{2}$ 为
\begin{equation}
E _ {\mathrm{R}} = E _ {0} - h \nu_ {\mathrm{e}} = \frac {E _ {0} ^ {2}}{2 m _ {\mathrm{Fe}} c ^ {2}} = 3. 1 \times 1 0 ^ {- 2 2} \mathrm{J}
\label{eq:7.s95}
\end{equation}
\end{subsolution}
\begin{subsolution}
设该原子核从静止的基态 $^{57}$Fe 吸收一频率为 $\nu_{a}$ 的光子，跃迁到激发态 $^{57}$Fe $^{*}$ ，并同时获得一反冲速度 $v'$ （吸收过程）。由能量、动量守恒有
\begin{equation}
h \nu_ {\mathrm{a}} = E _ {0} + \frac {1}{2} m _ {\mathrm{Fe}} v ^ {\prime 2}
\label{eq:8.s95}
\end{equation}
\begin{equation}
\frac {h \nu_ {\mathrm{a}}}{c} = m _ {\mathrm{Fe}} v ^ {\prime}
\label{eq:9.s95}
\end{equation}
由\eqref{eq:8.s95}\eqref{eq:9.s95}式解得
\begin{equation}
\nu_ {\mathrm{a}} = \frac {E _ {0}}{h} + \frac {h \nu_ {\mathrm{a}} ^ {2}}{2 m _ {\mathrm{Fe}} c ^ {2}} \approx \nu_ {0} + \nu_ {0} \frac {E _ {0}}{2 m _ {\mathrm{Fe}} c ^ {2}}
\label{eq:10.s95}
\end{equation}
这里也利用了 $\frac{h\nu_{a}}{m_{Fe}c^{2}}<<1$ ，将上式右边的 $\nu_{a}$ 用 $\nu_{0}$ 代替。于是
\begin{equation}
\nu_ {\mathrm{a}} - \nu_ {0} = \nu_ {0} \frac {E _ {0}}{2 m _ {\mathrm{Fe}} c ^ {2}} = 4. 7 \times 1 0 ^ {1 1} \mathrm{Hz}
\label{eq:11.s95}
\end{equation}
由\eqref{eq:8.s95}\eqref{eq:11.s95}式得，反冲能 $E_{R}^{\prime}=\frac{1}{2}m_{Fe}v^{\prime2}$ 为
\begin{equation}
E _ {\mathrm{R}} ^ {\prime} = h \nu_ {\mathrm{a}} - E _ {0} = \frac {E _ {0} ^ {2}}{2 m _ {\mathrm{Fe}} c ^ {2}} = E _ {\mathrm{R}} = 3. 1 \times 1 0 ^ {- 2 2} \mathrm{J}
\label{eq:12.s95}
\end{equation}
\end{subsolution}
\begin{subsolution}
由\eqref{eq:6.s95}\eqref{eq:11.s95}式有
\begin{equation}
\nu_ {\mathrm{e}} <   \nu_ {0} <   \nu_ {\mathrm{a}}
\label{eq:13.s95}
\end{equation}
由题设条件知，处于第一激发态 $^{57}Fe^{*}$ 能级的宽度为 $\Gamma$ 。静止的处于第一激发态原子核 $^{57}Fe^{*}$ 跃迁到基态时发出的光子的频率 $\nu$ 处于频率范围
\begin{equation}
\left[ \nu_ {\mathrm{e}} - \frac {\Gamma}{h}, \nu_ {\mathrm{e}} + \frac {\Gamma}{h} \right]
\label{eq:14.s95}
\end{equation}
内，能够被另一个静止的基态原子核 $^{57}$Fe 吸收而跃迁到第一激发态 $^{57}$Fe $^{*}$ 的光子的频率 $\nu'$ 处于频率范围
\begin{equation}
\left[ \nu_ {\mathrm{a}} - \frac {\Gamma}{h}, \nu_ {\mathrm{a}} + \frac {\Gamma}{h} \right]
\label{eq:15.s95}
\end{equation}
内。共振吸收能够发生的条件是频率范围\eqref{eq:14.s95}和\eqref{eq:15.s95}有重叠。频率范围\eqref{eq:14.s95}和\eqref{eq:15.s95}有重叠时应有
\begin{equation}
\nu_ {\mathrm{e}} - \frac {\Gamma}{h} <   \nu_ {\mathrm{a}} - \frac {\Gamma}{h} \leq \nu_ {\mathrm{e}} + \frac {\Gamma}{h} <   \nu_ {\mathrm{a}} + \frac {\Gamma}{h}
\label{eq:16.s95}
\end{equation}
考虑到\eqref{eq:13.s95}式可知，频率范围\eqref{eq:14.s95}和\eqref{eq:15.s95}有重叠的条件（即共振吸收能够发生的条件）是
\begin{equation}
\nu_ {\mathrm{a}} - \frac {\Gamma}{h} \leq \nu_ {\mathrm{e}} + \frac {\Gamma}{h}
\label{eq:17.s95}
\end{equation}
由\eqref{eq:7.s95}\eqref{eq:12.s95}\eqref{eq:17.s95}式得，\eqref{eq:17.s95}式等价于
\begin{equation}
\Gamma > E _ {\mathrm{R}}
\label{eq:18.s95}
\end{equation}
这与题给数据矛盾。实际上
\begin{equation}
\frac {E _ {\mathrm{R}}}{\Gamma} = \frac {\frac {E _ {0}}{2 m _ {\mathrm{Fe}} c ^ {2}} E _ {0}}{3 . 2 \times 1 0 ^ {- 1 3} E _ {0}} = 4. 2 \times 1 0 ^ {5}
\label{eq:19.s95}
\end{equation}
即
\begin{equation}
E _ {\mathrm{R}} > > \Gamma \quad (\text {大了} 5 \text {个数量级})
\label{eq:20.s95}
\end{equation}
因此不可能发生共振吸收。
\end{subsolution}
\begin{subsolution}
当原子核处于晶体中时，发射和吸收光子时晶体作为一个整体反冲，但因晶体的质量远大于原子核的质量，其反冲能可视为零。考虑到激发能级具有一定的宽度，发射的光子和能够吸收的光子的能量都应在
\begin{equation}
\left[ E _ {0} - \Gamma , E _ {0} + \Gamma \right]
\label{eq:21.s95}
\end{equation}
范围内，或光子频率应在上下限
\begin{equation}
\left[ \nu_ {0} - \frac {\Gamma}{h}, \nu_ {0} + \frac {\Gamma}{h} \right]
\label{eq:22.s95}
\end{equation}
之内的范围内。设晶体 A 相对于另一静止的晶体 B 以速度 $V$ （$V$ 的正负对应于两块晶体相互接近或相互远离）运动时，当 A 发射频率为 $\nu$ （在相对于 A 静止的参考系中）的光子，由于多普勒效应，B 所接收到的该光子的频率 $\nu'$ 为
\begin{equation}
\nu^ {\prime} = \nu \sqrt {\frac {1 + \beta}{1 - \beta}} \approx \nu (1 + \beta) = \left(1 + \frac {V}{c}\right) \nu
\label{eq:23.s95}
\end{equation}
式中 $\beta = \frac{V}{c}$ 。由此得
\begin{equation}
V = \left(\frac {\nu ^ {\prime}}{\nu} - 1\right) c = \left(\frac {E ^ {\prime}}{E} - 1\right) c
\label{eq:24.s95}
\end{equation}
式中 $E' = h\nu'$ ， $E = h\nu$ 。

在 A 与 B 相互接近的情形下，B 接收到的光子频率增加，当 A 发射的最低频率光子（对应于 $E = E_{0} - \Gamma$ ）的频率因多普勒效应增加到超过最高频率（对应于 $E = E_{0} + \Gamma$ ）时，共振吸收条件便无法满足。由\eqref{eq:24.s95}式得，V 的上限为
\begin{equation}
V _ {\max} = \left(\frac {E _ {0} + \Gamma}{E _ {0} - \Gamma} - 1\right) c
\label{eq:25.s95}
\end{equation}
在 A 与 B 相互远离的情形下，B 接收到的光子频率减小，当 A 发射的最高频率光子（对应于 $E = E_{0} + \Gamma$ ）的频率因多普勒效应减小到低于最低频率（对应于 $E = E_{0} - \Gamma$ ）时，共振吸收条件便无法满足。由\eqref{eq:24.s95}式得，V 的下限为
\begin{equation}
V _ {\min} = \left(\frac {E _ {0} - \Gamma}{E _ {0} + \Gamma} - 1\right) c
\label{eq:26.s95}
\end{equation}
因此，为使共振吸收发生，V 应处于下述范围内
\begin{equation}
- 0. 1 9 \mathrm{mm} \cdot \mathrm{s} ^ {- 1} = - \frac {2 \Gamma}{E _ {0}} c \approx - \frac {2 \Gamma}{E _ {0} + \Gamma} c \leq V \leq \frac {2 \Gamma}{E _ {0} - \Gamma} c \approx \frac {2 \Gamma}{E _ {0}} c = 0. 1 9 \mathrm{mm} \cdot \mathrm{s} ^ {- 1}
\label{eq:27.s95}
\end{equation}
上述分析表明，通过使吸收晶体以一定速度运动可补偿待测光子频率的变化；当速度合适时，恰好可实现共振吸收。只要不断地改变速度并测量吸收，就能通过达到共振吸收时速度的改变量来判断频率的变化，从而达到测量光子频率微小变化的目的。假设由于某种原因，到达吸收晶体处的光子频率 $\nu$ 相对于 $\nu_{0}$ 发生了相对变化 $\frac{\nu - \nu_{0}}{\nu_{0}} = \pm 1 \times 10^{-10}$ ，可让吸收晶体以一定的速度 $V'$ （ $V'$ 的正负表示晶体速度方向与光子运行方向相反或相同）运动，以补偿光子频率的上述变化，从而实现共振吸收，以精确测量光子频率的微小变化。为了使待测光子的频率 $\nu$ 在吸收晶体自身参考系中变为 $\nu_{0}$ ，由多普勒效应公式\eqref{eq:24.s95}，所需的吸收晶体的运动速度为
\begin{equation}
V ^ {\prime} = - \frac {\nu - \nu _ {0}}{\nu} c \approx - \frac {\nu - \nu _ {0}}{\nu _ {0}} c = \mp 1 \times 1 0 ^ {- 1 0} c = \mp 3 \mathrm{cm} \cdot \mathrm{s} ^ {- 1}
\label{eq:28.s95}
\end{equation}
通过测量达到共振吸收时吸收晶体的速度，就能够测出光子频率的改变量。
\end{subsolution}
\end{solution}